\section{Introduction}
\label{sec:snarkprobe-introduction}

We have seen a growing interest in privacy and privacy-enhancing technologies from the general public~\cite{adoption}, which has led to subsequent increased interest from parties that build and deploy the technologies that we use every day, with even the US government expressing interest in ways to best deploy privacy-enhancing technologies for data analytics~\cite{usgov}.
One of the key components in many privacy-enhancing protocols are zero-knowledge proofs~\cite{gmw86}, which are cryptographic algorithms that allow one party (the prover) to prove to another (the verifier) that a statement is true, without revealing any information beyond the validity of the statement itself.
Of these, one of the most popular instantiations are zkSNARKs (Zero-Knowledge Succinct Non-Interactive Arguments of Knowledge)~\cite{cryptoeprint:2013:507,parno2013pinocchio,groth2016size}, because of their small proof size, fast verification, and expressiveness.

Because of this popularity, we have seen an explosion of new zkSNARK protocols and libraries being developed in academia, and substantial interest from industry and other domains in actually deploying these protocols for real world usage~\cite{zcash,darkforest,aleo,filecoin,zkrollup}.


Unfortunately, they can be quite difficult to implement correctly, as the protocols themselves can be complicated and involved, and much of the work to generate a single proof for a single application is often done manually and hand-tuned to ensure maximum performance, with developers often working at the circuit or gate level to design the best protocols.  Additionally, as we have seen in the past, deploying complex cryptography has not come without its challenges~\cite{garman2016dancing,naveed2015inference,curveball,rupprecht2020call}, and applications that use zkSNARKs have not been immune to these either, as Zcash for example has found various cryptographic bugs in different components~\cite{zcashcounterfeit,zcashvulns,zcashaudit}.  And checking any cryptographic implementation manually can be time consuming and potentially error prone.

While automated techniques like fuzzing have been used to test cryptographic implementations specifically before~\cite{aumasson2017automated}, we cannot use existing cryptographic fuzzers here as they generally test only for known weaknesses in certain schemes and are unable to produce the types of inputs that we need in a cost-effective way.
Additionally, most general purpose fuzzing tools cannot detect what we refer to as \textit{\logicerror{s}}, i.e., errors that might result in an incorrect computation but not a program crash, particularly with regards to \zksnark libraries.

All of this motivates us to ask the question:

\smallskip
\noindent\emph{Can we develop better tooling to automatically check the security of and precisely locate software bugs and cryptographic logic errors in the proof generation processes and libraries of zkSNARK protocols and the applications that use them?}
\smallskip

To help answer this question in the affirmative, we design and build \tool to automatically both check the correctness and consistency of proof programs generated by R1CS-based zkSNARK libraries, as well as inspect the security of the libraries themselves and flag any inconsistencies between a protocol's implementation and its description, no matter how minor\footnote{Note that this is something that professional security audits do flag, as these can lead to potential future problems, see~\cite{zcashaudit}.}.  Our primary goal is to make the process as automated as possible, thus reducing the chance of human error in the process, as well as lowering the barrier to entry for usage.  We wish to enable even those who might not be cryptographic experts to inspect a library or proof implementation without understanding details of the underlying protocol or specifics of the library.  Users need only indicate some configuration settings such as the fuzzer parameters and expected proof statement, and our tool will automatically analyze the library and proof program to detect possible implementation and cryptography related errors.

In order to achieve this, we utilize a combination of techniques.
Dynamic analysis allows us to trace real-time data and variable values, as well as handle a variety of different zkSNARK libraries written in different programming languages without needing additional (manual) language specific adaptations.
Custom fuzzing techniques allow us to produce a variety of valid or invalid R1CS matrices (i.e., proofs) to exercise the different codepaths in a library.
SMT (Satisfiability Modulo Theory) solvers allow us to help verify the consistency of a user-specified set of statement equations (i.e., the desired proof) with the actual R1CS matrix used in the application.
And the notion of ideal files and a value checker helps ensure that a library correctly realizes the given underlying protocol.

While one could tackle some of this from a formal analysis approach~\cite{projecteverest}, we view \tool to be complementary, as formal analysis works best for specific, static protocol implementations (like libraries), and we wish to also be able to check application-specific proofs.  Formal verification in such a scenario would need to be done for \textit{each individual proof}, which is both time-consuming and requires expertise in such techniques for every project that wishes to use a zkSNARK. 
\tool, on the other hand, is designed to work with a variety of different protocols and to be easy to run on a given proof implementation in an application, without requiring substantial expertise. 
We will see later on how a formally verified library for a specific protocol could be adapted as part of \tool to provide stronger guarantees, though we leave this for future work.
As such, we believe \tool is a step in the right direction for ensuring the correctness of existing zkSNARK implementations and applications.

We choose to focus on two widely used zkSNARK protocols, \pinocchio~\cite{parno2013pinocchio,ben2014succinct} and \groth16~\cite{groth2016size}, and four libraries, \libsnark~\cite{libsnark}, \bellman~\cite{bellman}, \arkworks~\cite{arkworks}, and \gnark~\cite{gnark}. \libsnark implements both \pinocchio and \groth16 in C++, while \bellman and \arkworks implement \groth16 in Rust, and \gnark implements \groth16 in Go.  These represent four of the most popular and widespread open source libraries that have seen real world usage and underpin many projects that use zkSNARKs.  This selection also demonstrates the flexibility of \tool to handle different R1CS-based zkSNARK protocols, as well as diverse languages.

In addition to designing and implementing \tool, we tested its performance extensively, using a number of different statement types to test the \ConstraintChecker, and a wide variety of R1CS matrix sizes and configurations (which are representative of statements of varying complexity) to test \SnarkFuzzer.

We demonstrate performance that we believe is both reasonable and scalable for existing zkSNARK use cases\footnote{And, in fact, in one instance the underlying library was actually unable to scale and run before we ran into any issues with our tool.}.
We also evaluated \tool's ability to detect potential inconsistencies or different types of errors.
While we did not find any new exploitable errors, we did find a number of issues across the various libraries that we consider to be potentially unsafe or that might require special care by a developer to navigate correctly.

% \parhead{Our contributions.} In this paper we make the following contributions:
\subsection{Our contributions}
In this paper we make the following contributions:

%\vspace*{-5pt}
\begin{itemize}[leftmargin=*]
    \setlength{\itemsep}{0pt}
    \setlength{\parskip}{0pt}
		\item We design \tool, a tool that can automatically check for potential software errors and cryptographic logic errors, as well as for consistency in the implemented proof statement, in R1CS-based zkSNARK libraries.
		\item We implement and build \tool\footnote{https://github.com/BARC-Purdue/SNARKProbe} for two zkSNARK protocols, \pinocchio~\cite{parno2013pinocchio,ben2014succinct} and \groth16~\cite{groth2016size}, and four libraries, \libsnark~\cite{libsnark}, \bellman~\cite{bellman}, \arkworks~\cite{arkworks}, and \gnark~\cite{gnark}, to demonstrate its flexibility and extensibility.
		\item To the best of our knowledge, we are the first to explore applying fuzzing and other analysis techniques to zkSNARKs specifically.
	  \item We evaluate the performance of \tool in an extensive set of experiments, demonstrating its acceptable performance for real world applications.
		\item We demonstrate and discuss \tool's ability to \emph{automatically} catch different types of errors in \zksnark libraries, finding seven inconsistencies and successfully locating a prior CVE.
\end{itemize}

\subsection{Chapter Organization}
In Section~\ref{sec:snarkprobe-background}, we discuss the necessary background and related work.  In Section~\ref{sec:snarkprobe-overview}, we provide a brief overview of \tool, as well as its intended use cases.  In Section~\ref{sec:snarkprobe-design}, we introduce the high level design and goals for each component of our tool with the concrete tools and techniques that allow us to realize these design goals.  In Section~\ref{sec:snarkprobe-evaluation}, we evaluate the performance of \tool, and in Section~\ref{sec:snarkprobe-errorcatching} we discuss its ability to detect known errors.  Finally, we conclude in Section~\ref{sec:snarkprobe-conclusion}.

\vspace{1em}
\parhead{Possible ethical considerations.}
While our tool did find a few inconsistencies in gadgets and protocol specification versus implementation throughout our testing, all of these issues were either not exploitable, already known, or just require care from a developer when implementing a proof, and as such we do not believe there is anything to disclose.
